\documentclass[11pt]{article}
\usepackage[margin=1in]{geometry}
\usepackage{hyperref}
\usepackage{listings}
\usepackage{xcolor}
\usepackage{booktabs}
\usepackage{enumitem}
\usepackage{longtable}
\usepackage{array}
\usepackage{amsmath}

\lstset{
  basicstyle=\ttfamily\small,
  breaklines=true,
  columns=fullflexible,
  frame=single,
  backgroundcolor=\color{gray!10},
  showstringspaces=false
}

\title{Exploratory Analysis of the URI OPeNDAP Server\\
       \large \texttt{https://sst-aqua.gso.uri.edu/opendap/}\\
       \large grep-dap test case, Prompt \#1}
\author{Compiled for P. Cornillon}
\date{2026-05-16}

\begin{document}
\maketitle

\begin{abstract}
This document records what an outside client can learn about the URI
OPeNDAP server at
\url{https://sst-aqua.gso.uri.edu/opendap/} given only its
catalog pages, the DAP4 metadata documents (DMRs) of the publicly
readable datasets, and small subsampled slices of the actual data.
The server is a Hyrax instance hosting 17 top-level containers, of
which only \texttt{SST\_Orbits/} and \texttt{gradients\_by\_period/}
are publicly readable; the rest return Hyrax 403 (BES filesystem
permissions). The two readable branches sit at opposite ends of the
semantic-metadata spectrum that motivates \texttt{grep-dap}: the
orbit files are richly CF-annotated, and the gradient-statistics
files carry essentially no semantic metadata at all. The report
describes each branch, distinguishes what is stated by the metadata
from what is inferred by structure or by reading the data, and notes
specific places where a metadata-inference tool would have to do
useful work.
\end{abstract}

\tableofcontents

\section{Server identification}

\begin{description}[leftmargin=2em,style=nextline,topsep=2pt]
  \item[Software.] Hyrax (OPeNDAP, Inc.), behind \texttt{nginx}.
    Confirmed by (i) the \texttt{Server: nginx} header on the
    landing page, (ii) the rendered \texttt{contents.html}
    template that includes the OPeNDAP logo, the
    \texttt{contents.css} stylesheet, and the standard ``DAP4
    Response Links'' / ``DAP2 Response Links'' columns, and
    (iii) the Hyrax-format 403 error page on restricted
    directories, which is BES-style and names the on-disk path
    (\texttt{/usr/share/hyrax/MUR/} and similar).
  \item[Protocol.] Both DAP2 and DAP4 are advertised on every
    dataset; DAP4 is the natural choice for this project because
    a single \texttt{.dmr.xml} fetch returns structure plus all
    attributes. \texttt{pydap.client.open\_url(\dots,
    protocol="dap4")} succeeds on both readable branches.
  \item[Authentication.] No challenge was issued. The 403s are
    BES filesystem permissions, not \texttt{401 Unauthorized}; a
    typical Earthdata-style login is not in play. Whatever
    credentials would unlock the other 15 directories live on the
    server itself.
  \item[Hostname signal.] The hostname \texttt{sst-aqua.gso.uri.edu}
    nominates the URI Graduate School of Oceanography
    (\texttt{gso.uri.edu}) and points at MODIS-Aqua SST as the
    likely subject. Both expectations are confirmed by the
    catalog contents (see \S\ref{sec:sstorbits}).
\end{description}

\section{Top-level catalog}\label{sec:catalog}

The root carries 17 entries (one of which, \texttt{@eaDir/}, is a
Synology NAS housekeeping directory and is ignored). Of the remaining
16, two are publicly readable and 14 return Hyrax 403.

\begin{center}
\begin{tabular}{@{}lll@{}}
\toprule
Container & Access & Likely contents (from name alone) \\
\midrule
\texttt{SST\_Orbits/}             & \textbf{200} & MODIS-Aqua L2 SST orbit files (this study) \\
\texttt{gradients\_by\_period/}   & \textbf{200} & Monthly SST + gradient statistics (this study) \\
\midrule
\texttt{JAXA\_Orbits/}            & 403 & JAXA satellite (likely AMSR2 / GCOM-W) orbit data \\
\texttt{JAXA\_Orbits\_OLD/}       & 403 & Older JAXA orbit holdings \\
\texttt{MUR/}                     & 403 & JPL MUR L4 SST analysis \\
\texttt{RSS\_Orbits/}             & 403 & Remote Sensing Systems microwave SST orbits \\
\texttt{RSS\_Orbits\_OLD/}        & 403 & Older RSS holdings \\
\texttt{RSS\_Orbits\_REALLY\_OLD/}& 403 & Still-older RSS holdings \\
\texttt{gradients\_by\_period\_5day/} & 403 & 5-day analogue of \texttt{gradients\_by\_period} \\
\texttt{iQuamBuoy/}               & 403 & NOAA iQuam in-situ SST (drifter / buoy) \\
\texttt{matchups/}                & 403 & Satellite$\leftrightarrow$in-situ matchup tables \\
\texttt{matchups\_v2/}            & 403 & Successor matchup version \\
\texttt{matchups\_v3/}            & 403 & Successor matchup version \\
\texttt{timeSeries/}              & 403 & Per-location SST time series \\
\texttt{timeSeries\_v2/}          & 403 & Successor time-series version \\
\texttt{timeSeries\_v2\_backup/}  & 403 & Backup snapshot of v2 \\
\bottomrule
\end{tabular}
\end{center}

The 14 names map cleanly onto four thematic axes that any inferring
client should record even when it cannot read the contents:
\begin{itemize}
  \item \textbf{Sensor / agency} (JAXA, RSS, MUR, AQUA implicit in
    the hostname).
  \item \textbf{Geometry} (Orbits = L2 swath; gradients-by-period =
    gridded statistics; timeSeries = per-location).
  \item \textbf{Versioning} (\texttt{\_v2}, \texttt{\_v3},
    \texttt{\_OLD}, \texttt{\_REALLY\_OLD}) --- a process history.
  \item \textbf{In-situ vs.\ satellite vs.\ matchup}
    (\texttt{iQuamBuoy} is in-situ, \texttt{matchups[\_v2,\_v3]}
    are by construction satellite\,$\times$\,in-situ collocations).
\end{itemize}

That is meaningful catalog-level metadata extracted from
\emph{nothing but folder names}, and it is the kind of inference
\texttt{grep-dap} should preserve.

\section{Readable branch \#1: \texttt{SST\_Orbits/}}\label{sec:sstorbits}

\subsection{Directory structure and filenames}

The branch is partitioned by calendar year (\texttt{2002}--\texttt{2024})
and within each year by month (\texttt{01}--\texttt{12}). Leaf
filenames carry a great deal of structured information:

\begin{lstlisting}
AQUA_MODIS_orbit_115220_20240101T011558_L2_SST-URI_24-2.nc4
^^^^^^^^^^      ^^^^^^^ ^^^^^^^^^^^^^^^ ^^^^^^^^^^^ ^^^^
sensor          orbit#  start (ISO)     product     version
\end{lstlisting}

So even before opening a single file we can write down:
\begin{itemize}
  \item Source platform: NASA Aqua, MODIS instrument.
  \item Processing level: L2 (per-orbit, swath-projected, no global
    composite).
  \item Product family: SST.
  \item Producer / processor: URI (the host's own institution).
  \item Local version: \texttt{24-2}. Coupled with the
    \texttt{\_OLD} naming pattern of sibling directories this is
    consistent with URI maintaining its own reprocessing series.
  \item Temporal coverage: 2002-07-04 (MODIS-Aqua first light is
    early July 2002) onward. The earliest filename in the
    catalog will pin this exactly; the present sample is in
    January 2024.
  \item Volume: $\sim$14.4 orbits per day $\times$ 365\,d $\times$
    23 years $\approx 1.2\times 10^5$ orbit files.
\end{itemize}

\subsection{Structural metadata (DMR)}

The DMR for the sample orbit
\begin{lstlisting}
https://sst-aqua.gso.uri.edu/opendap/SST_Orbits/2024/01/
    AQUA_MODIS_orbit_115220_20240101T011558_L2_SST-URI_24-2.nc4
\end{lstlisting}
declares three root-level dimensions and two groups (DAP4-only):

\begin{center}
\begin{tabular}{@{}lll@{}}
\toprule
Dimension & Size & Interpretation \\
\midrule
\texttt{ny} & 40\,271 & along-track scan lines of the orbit \\
\texttt{nx} & 1\,354  & across-track pixels (MODIS 1\,km swath width) \\
\texttt{i}  & 5       & along-orbit subregion index (5 segments) \\
\bottomrule
\end{tabular}
\end{center}

Root-level variables (selected; types and shapes drawn directly from
the DMR; semantic content from the DMR attributes):

\begin{center}
\renewcommand{\arraystretch}{1.1}
\begin{tabular}{@{}lp{0.20\linewidth}p{0.16\linewidth}p{0.42\linewidth}@{}}
\toprule
Type & Name & Shape (dims) & Notes from DMR attributes \\
\midrule
Int16 & \texttt{SST\_In}                              & (ny, nx) & \texttt{units=C}, scale 0.005, \texttt{standard\_name=sea\_surface\_temperature}, valid\_min/max $-600/9000$ raw. \\
Int16 & \texttt{regridded\_sst}                       & (ny, nx) & \texttt{units=C}, scale 0.005, \texttt{standard\_name=masked\_sea\_surface\_temperature}. A smoothed / masked SST field. \\
Int32 & \texttt{eastward\_gradient}                   & (ny, nx) & \texttt{units=C/km}, scale 0.0001, \texttt{standard\_name=eastward\_temperature\_gradient}. \\
Int32 & \texttt{northward\_gradient}                  & (ny, nx) & \texttt{units=C/km}, scale 0.0001, \texttt{standard\_name=northward\_temperature\_gradient}. \\
Int8  & \texttt{refined\_mask}                        & (ny, nx) & Mask removing original high-gradient pixels (long\_name). \\
Int8  & \texttt{qual\_sst}                            & (ny, nx) & Per-pixel SST quality flag. \\
Int32 & \texttt{latitude}, \texttt{longitude}         & (ny, nx) & Pixel coordinates. \texttt{units=degrees\_north/east} but no scale\_factor declared; raw values are $10^3 \times$ degrees (see \S\ref{sec:gaps}). \\
Int32 & \texttt{regridded\_latitude}, \texttt{regridded\_longitude} & (ny, nx) & Same as above but on the regridded grid. \\
Float32 & \texttt{nadir\_latitude/longitude}, \texttt{left/right\_swath\_edge\_trackline\_latitude/longitude} & (ny,) & Swath geometry, in degrees. \\
Float32 & \texttt{time\_from\_start\_orbit}          & (ny,) & Seconds since orbit start. \\
Float64 & \texttt{DateTime}                          & scalar  & Seconds since 1970-01-01 UTC; \texttt{standard\_name=time}. \\
Int32   & \texttt{region\_start}, \texttt{region\_end} & (i=5,) & Index ranges of the five along-orbit subregions. \\
\bottomrule
\end{tabular}
\end{center}

\paragraph{Group \texttt{/Regrid\_to\_L2eqa/}.} The orbit file is not
just an L2 SST file: it carries an additional group of fields
\emph{regridded to an L2 equal-area grid}, together with collocated
\textbf{AMSR-E} (passive microwave) wind speed, water vapor, cloud
liquid water, and SST. Notable variables:
\texttt{L2eqa\_MODIS\_SST}, \texttt{L2eqa\_MODIS\_std\_SST},
\texttt{L2eqa\_MODIS\_num\_SST} (MODIS resampled with sub-cell
statistics), \texttt{L2eqa\_AMSR\_E\_SST/wind\_speed/water\_vapor/
cloud\_liquid\_water} (AMSR-E resampled), the native
\texttt{AMSR\_E\_lat/lon/SST/\dots}, plus
\texttt{MODIS\_SST\_on\_AMSR\_E\_grid} and
\texttt{convolved\_MODIS\_on\_AMSR\_E\_grid}. This is the
infrastructure of a MODIS$\leftrightarrow$AMSR-E intercomparison
product.

AMSR-E stopped producing data in October 2011, so the AMSR-E-named
variables in 2024 orbit files almost certainly contain only fill
values --- the schema persists past the instrument's lifetime. A
2010-era file would be the place to check this.

\paragraph{Group \texttt{/contributing\_granules/}.} Provenance:
filenames, start/end times, and byte-range indices of the source
NASA L2 granules that were stitched together to form the orbit. No
science, but valuable for reproducibility tracking.

\subsection{Sampled data}

Two slice reads with \texttt{pydap} validate the metadata; see
\texttt{scripts/uri\_sample\_data.py},
\texttt{scripts/uri\_sample\_equator.py}.

\paragraph{Polar slice} ($i_0=20000$, $j_0=600$, $100\times 100$).
Nadir is at $\sim 78\text{--}79^\circ$N, longitude
$\sim 45\text{--}49^\circ$E (Kara Sea region). \texttt{SST\_In}
ranges $-15$ to $-3\,^\circ$C in scaled units; values below
$-3\,^\circ$C lie outside the declared \texttt{valid\_min/valid\_max},
consistent with sea ice / land / unconverged retrievals.
\texttt{eastward\_gradient} and \texttt{northward\_gradient} return
empty (every pixel is \texttt{\_FillValue}) at this latitude.

\paragraph{Equatorial slice} ($i_0 \approx 29259$, $j_0=600$,
$100\times 100$). Nadir at $\sim 0^\circ$N, $\sim 7^\circ$W (off the
African coast). \texttt{SST\_In} in $[16.2, 29.0]\,^\circ$C, mean
$27.4$, std $1.7$. \texttt{eastward\_gradient},
\texttt{northward\_gradient} populate at $\sim$20\,\% of pixels
(1\,993 of 10\,000) with magnitudes mostly within $\pm
0.5\,^\circ$C/km, mean near zero --- as expected for open tropical
ocean with occasional cloud edges. \texttt{regridded\_sst} populates
at 36\,\%, range $[26.8, 29.0]\,^\circ$C --- a smoothed / more
heavily masked version of the same field.

\paragraph{Time-stamp self-consistency.} The scalar
\texttt{DateTime} returns $1.7041 \times 10^9$\,s, which is
2024-01-01T01:15:58 UTC, matching the filename's
\texttt{20240101T011558} token exactly.

\subsection{Metadata gaps discovered in sampling}\label{sec:gaps}

\begin{itemize}
  \item The Int32 \texttt{latitude} / \texttt{longitude} arrays
    declare \texttt{units=degrees\_north} / \texttt{degrees\_east}
    but do \emph{not} declare a \texttt{scale\_factor}. Raw values
    sampled near the equator are $\sim\pm 500$ (i.e.\ $\pm 0.5$
    after dividing by $10^3$), and near the pole are
    $\sim 7.9 \times 10^4$ (i.e.\ $79^\circ$ after dividing by
    $10^3$). Verified against the Float32 \texttt{nadir\_latitude}
    and \texttt{nadir\_longitude} (which are in degrees, plain).
    \emph{The implicit scale on the Int32 coordinates is $10^{-3}$;
    this is structural information present nowhere in the DMR.}
  \item No global (root-group) attributes are exposed in the DMR.
    No \texttt{Conventions}, no \texttt{title}, no
    \texttt{institution}, no \texttt{processing\_level}, no
    \texttt{platform}, no \texttt{instrument} attribute. All of
    that has to be inferred from the filename and the variable
    naming.
\end{itemize}

\section{Readable branch \#2: \texttt{gradients\_by\_period/}}\label{sec:gradients}

\subsection{Layout}

A flat directory of 270 NetCDF files,
\texttt{augmented\_monthly\_stats\_MM\_YYYY.nc}, spanning
\texttt{01\_2003} through \texttt{12\_2024}. ($22 \times 12 = 264$
with a small surplus that I have not chased down --- likely a small
number of stray files or revision tags.) The filename pattern is
the only temporal metadata: there is no explicit
\texttt{time} variable inside.

\subsection{Structural metadata (DMR)}

\begin{center}
\begin{tabular}{@{}lll@{}}
\toprule
Dimension & Size & Inferred meaning \\
\midrule
\texttt{lon} & 360 & $1^\circ$ longitude bins \\
\texttt{lat} & 180 & $1^\circ$ latitude bins \\
\bottomrule
\end{tabular}
\end{center}

34 variables, all \texttt{Float64} of shape \texttt{(lon, lat)}.
\textbf{No semantic metadata.} No \texttt{units}, no \texttt{long\_name},
no \texttt{standard\_name}, no \texttt{\_FillValue}, no
coordinate variables, no global attributes. This is the OPeNDAP
``free-for-all'' that motivates \texttt{grep-dap}.

\subsection{What the names imply}

The variable names follow a consistent grammar:

\begin{lstlisting}
{day|night}_{sum|}_{quantity}[_squared]
{day|night}_{pixel|as_pixel|at_pixel}_count
\end{lstlisting}

That decomposes cleanly into:

\begin{description}[leftmargin=2em,topsep=2pt]
  \item[Period partition.] \texttt{day\_} vs.\ \texttt{night\_} ---
    presumably the MODIS Aqua $\sim$13:30 local ascending vs.\
    $\sim$01:30 descending nodes, partitioned at the per-pixel
    level upstream.
  \item[Bookkeeping.] For each quantity we get a sum
    ($\Sigma$) and a sum of squares ($\Sigma^2$), the canonical
    pair for computing mean and variance after aggregation. The
    associated pixel-count $N$ closes the loop.
  \item[Quantity flavors.] Six families of derived statistics:
    \begin{enumerate}
      \item \texttt{SST} (the scalar SST values themselves).
      \item \texttt{eastward\_gradient} (east component, signed).
      \item \texttt{northward\_gradient} (north component, signed).
      \item \texttt{magnitude\_gradient} (per-pixel magnitude in
        unspecified units; consistent with the orbit-file
        magnitude having been summed in geographic coordinates).
      \item \texttt{grad\_as\_per\_km} (along-scan component, per
        km).
      \item \texttt{grad\_at\_per\_km} (along-track component, per
        km).
      \item \texttt{grad\_mag\_per\_km} (per-km magnitude).
    \end{enumerate}
    (Six conceptual families, but $4+3$ named gradient variables
    because magnitude appears in both swath-relative and
    geographic coordinate triples.)
  \item[Three count flavors.] \texttt{pixel\_count} (valid SST),
    \texttt{as\_pixel\_count} (valid along-scan gradient), and
    \texttt{at\_pixel\_count} (valid along-track gradient). The
    fact that the three counts differ confirms that the gradient
    fields are masked more aggressively than the SST field
    itself, which matches the orbit-file observation that gradient
    pixels were entirely empty in the polar slice while SST
    pixels were not.
\end{description}

\subsection{Sampled data}\label{sec:gradients-samples}

For \texttt{augmented\_monthly\_stats\_07\_2020.nc}, full $360\times
180$ reads via \texttt{pydap}:

\begin{center}
\begin{tabular}{@{}lrr@{}}
\toprule
Variable & valid $n$ & range (min, max) \\
\midrule
\texttt{day\_pixel\_count}              & 64\,800 & $0$, \quad $2.02 \times 10^{5}$ \\
\texttt{night\_pixel\_count}            & 64\,800 & $0$, \quad $1.96 \times 10^{5}$ \\
\texttt{day\_sum\_SST}                  & 63\,431 & $-628$, $5.28 \times 10^{6}$ \\
\texttt{night\_sum\_SST}                & 64\,028 & $-947$, $5.44 \times 10^{6}$ \\
\texttt{day\_sum\_magnitude\_gradient}  & 64\,800 & $0$, $1.37 \times 10^{4}$ \\
\texttt{day\_sum\_eastward\_gradient}   & 64\,800 & $-4\,483$, $2\,587$ \\
\texttt{day\_sum\_northward\_gradient}  & 64\,800 & $-2\,901$, $6\,113$ \\
\bottomrule
\end{tabular}
\end{center}

Two derived quantities:
\begin{itemize}
  \item Per-cell mean SST $\Sigma_{\rm SST} / N_{\rm pix}$ produced
    NaNs in 24\,900 cells with non-zero $N_{\rm pix}$, indicating
    that \texttt{day\_sum\_SST} is itself NaN in those cells (NaN
    is acting as an undeclared fill value). The clean cells fall
    in the plausible $-2$ to $32\,^\circ$C ocean range
    (range printed in the script log; full distribution not yet
    plotted).
  \item Per-cell mean magnitude gradient $\Sigma_{|\nabla|} /
    N_{\rm as}$ has mean $0.083\,^\circ$C/km and a tail to
    $22\,^\circ$C/km. The bulk is consistent with typical ocean
    SST gradients of $0.01\text{--}0.5\,^\circ$C/km; the tail is
    almost certainly numerical (cloud-edge / coast-edge
    artifacts).
\end{itemize}

\subsection{Inferred semantics for grep-dap to record}

Even with zero attributes in the file, the following can be
reconstructed with high confidence:

\begin{center}
\begin{tabular}{@{}lll@{}}
\toprule
Variable family & Inferred units & Source of inference \\
\midrule
\texttt{*\_pixel\_count} / \texttt{*\_as\_pixel\_count} / \texttt{*\_at\_pixel\_count} & dimensionless (count) & Name + values are nonneg integers in Float64. \\
\texttt{*\_sum\_SST}, \texttt{*\_sum\_SST\_squared} & $^\circ$C, $^\circ$C${}^2$ & SST in the source orbit files is in $^\circ$C; derived sums inherit units. \\
\texttt{*\_sum\_eastward\_gradient}, \texttt{*\_sum\_northward\_gradient} & $^\circ$C/km & Source orbit gradients are in $^\circ$C/km. \\
\texttt{*\_sum\_magnitude\_gradient} & $^\circ$C/km & Same. \\
\texttt{*\_sum\_grad\_\{as,at,mag\}\_per\_km}, with \texttt{\_squared} & $^\circ$C/km, $(^\circ$C/km$)^2$ & Suffix \texttt{\_per\_km} states the unit; \texttt{as} / \texttt{at} are swath-relative components. \\
\bottomrule
\end{tabular}
\end{center}

What \emph{cannot} be inferred from the file alone (without
re-examining the orbit data the file is plausibly derived from):

\begin{itemize}
  \item The orientation of the \texttt{lat} axis (south-to-north
    or reverse) and the convention for the \texttt{lon} axis
    ($0\to 360$ or $-180\to +180$). The absence of coordinate
    variables makes this irrecoverable from metadata alone.
    \emph{Resolvable by sampling}: check, for example, whether
    the cell with maximum mean SST sits near the equator
    (lat-index $\approx 90$ either way) and use the longitude
    pattern of mean SST to locate the warm pool and the equatorial
    Atlantic cold tongue.
  \item The exact definition of \texttt{magnitude\_gradient}
    relative to the eastward/northward components --- whether the
    summed magnitude in a cell equals the magnitude of the summed
    components (almost certainly no) or the sum of per-pixel
    magnitudes (likely yes). Resolvable analytically only by
    accessing the upstream pixel-level data; \texttt{grep-dap}
    can flag the question but not answer it.
  \item Whether $\Sigma$ and $N$ count the same pixels (e.g.\
    whether a NaN'd $\Sigma_{\rm SST}$ in a cell with $N_{\rm
    pix}>0$ means ``$N$ counts gradients, $\Sigma$ counts SST''
    or means ``arithmetic overflow / undeclared fill''). The
    sample shows the discrepancy clearly; the cause requires
    asking the data provider.
\end{itemize}

\section{Cross-branch picture}\label{sec:crossbranch}

The two readable branches almost certainly form a single pipeline:
the orbit files in \texttt{SST\_Orbits/} are the L2 input from which
the monthly gridded statistics in \texttt{gradients\_by\_period/}
are computed. Three pieces of evidence support that:

\begin{enumerate}
  \item Vocabulary match. \texttt{eastward\_gradient},
    \texttt{northward\_gradient}, and \texttt{magnitude\_gradient}
    appear in both files with the same semantics; units in the
    orbit file (where stated) are \texttt{C/km}, which is what
    the gradient-file naming (\texttt{\_per\_km}) implies.
  \item Quantity match. The day/night, along-scan/along-track
    split in the gradient files mirrors the way MODIS Aqua data
    are partitioned (ascending vs.\ descending, swath-relative
    derivatives) and how the orbit file's masking variables are
    organized.
  \item Volume match. $14.4 \text{ orbits/day} \times 31 \text{ d}
    \approx 446$ orbits per month; each orbit has
    $40\,271 \times 1\,354 = 5.45 \times 10^7$ pixels; a per-cell
    \texttt{day\_pixel\_count} max of $2.02 \times 10^5$ is
    consistent with $\sim$half-orbit (day-only) coverage at MODIS
    1\,km resolution binned into $1^\circ$ cells near the equator
    (where one $1^\circ$ cell is roughly $111 \times 111$\,km,
    i.e.\ $\sim 1.2 \times 10^4$ pixels per pass; with $\sim$15
    overlapping day passes a month at the equator, $\sim 2 \times
    10^5$ is the expected upper bound).
\end{enumerate}

The 403 directories that I could not read complete the picture but
only hypothetically: \texttt{matchups[\_v2,\_v3]/} plus
\texttt{iQuamBuoy/} suggest a validation chain (satellite SST vs.\
buoy SST), and the JAXA/RSS/MUR directories together with
\texttt{timeSeries[\_v2]/} suggest a multi-source SST archive whose
top product family is precisely what URI / NASA SST researchers
typically build.

\section{What this means for grep-dap}\label{sec:meaning}

The exercise produced six concrete observations that bear on the
project goals stated in \texttt{CLAUDE.md} and
\texttt{docs/opendap\_readme.tex}.

\begin{enumerate}
  \item \textbf{Catalog walking is half the inference.} Just from
    folder and file names, \texttt{grep-dap} can extract platform,
    sensor, level, producer, version, and rough temporal extent
    --- before reading a single byte of data. Catalog-name parsing
    should be a first-class capability, not a fallback.
  \item \textbf{Inaccessible directories are still signal.} Names
    of 403'd directories partition by sensor source, geometry,
    versioning, and matchups even when their contents are
    unreadable. \texttt{grep-dap} should not drop unreadable
    entries; it should retain them as catalog facts.
  \item \textbf{The semantic-metadata distribution is bimodal on
    this server.} The orbit files are fully CF-flavored (units,
    standard\_name, scale\_factor, valid\_min, valid\_max,
    \_FillValue); the derived monthly files have nothing. Real
    OPeNDAP archives plausibly look like this --- a few rich
    files surrounded by many sparse ones. Inference strategies
    that work only when rich attributes are present will not be
    enough.
  \item \textbf{Even rich CF metadata has gaps.} The orbit
    \texttt{latitude/longitude} arrays declare units but omit
    \texttt{scale\_factor}, requiring a $10^3$ rescale that no
    naive client will apply. \texttt{grep-dap} should be alert to
    declared units that are internally inconsistent with raw
    magnitudes for the declared dtype.
  \item \textbf{Small reads anchor the inference cheaply.} A
    single $100\times 100$ slice at the equator placed the SST
    field firmly in the 16--29$\,^\circ$C ocean range and
    confirmed gradient units and lat/lon scales. The cost was
    one DAP4 hyperslab request per variable. For an inferring
    agent, two or three such reads per dataset is a very low
    price.
  \item \textbf{Names are not enough on their own.} The
    \texttt{magnitude\_gradient} / \texttt{grad\_mag\_per\_km}
    duality in the gradient files cannot be disambiguated from
    the schema; only by tracing the orbit-file recipe (or by
    asking the producer) can the two be reconciled. An honest
    inference tool flags such ambiguities rather than guessing.
\end{enumerate}

\section{Process and reproducibility}

All probes and samples are scripted; the scripts live in
\texttt{scripts/}:

\begin{lstlisting}
uri_root_walk.py           # parse the root contents.html
uri_walk_public.py         # walk the two publicly readable branches
uri_walk_sst_orbits.py     # SST_Orbits structure walk
uri_sample_data.py         # pydap polar slice + gradient stats
uri_sample_equator.py      # pydap equatorial slice
\end{lstlisting}

The chronological log of this exploration is in
\texttt{docs/test\_case\_uri\_log.tex} with timestamps.

\section{Caveats}

\begin{itemize}
  \item Two of seventeen directories were readable; the picture
    above is therefore necessarily incomplete.
  \item The single orbit and the single monthly-stats file picked
    for inspection are representative of their branches by
    construction (one is the first orbit of the most recent full
    year; the other is mid-2020), not by random sampling.
  \item No claim about the global attribute set on the
    \emph{inaccessible} directories is made; everything written
    about them is inferred from the directory names only.
  \item The hostname expectation (MODIS-Aqua SST) and the
    catalog contents happen to agree. \texttt{grep-dap} must
    record both as separate facts and never conflate ``it is
    plausible'' with ``it is verified.''
\end{itemize}

\end{document}
